\subsection{The complete conditional integral comparison in rank four}
\label{sec:rh-symplectic-length}

We use the actual RH model of
Theorem~\ref{thm:rh-symplectic-model}. Keep its ambient augmented
domain \(R\), quotient \(S=R/J_{\rm eq}\), coefficient DVR \(O\)
with fraction field \(E_0\), and two weights
\[
 w_i=(\lambda_i,\mu_i),\qquad
 \lambda=(1,0,1,0),\quad \mu=(1,1,0,0).
\]
The grading of matrix entry \(ij\) is \(w_i-w_j\).
The three nonzero first-coordinate root parameters in the rational
chart have weights \(A=(1,0), B=(-1,1), T=(1,1)\).
The other weights are \(0,P=(0,1),-P\), with multiplicities
\(m,r',r\), respectively. All equations of \(J_{\rm eq}\)
are the original lists \(F,E_+,Z\) of weights \(-P,P,0\),
with respective sizes \(r+1,r',m\); in particular \(r\ge1\).

Here are the three remaining arithmetic premises, stated at the
actual points of this ring:
\begin{enumerate}
\item[(F)] The full algebra \(S_{(0,0)}\) of weight-zero
elements is finite over \(O\).
\item[(E)] No characteristic-zero point of \(S\) that is
block diagonal for the coordinate blocks \(\{0,2\},\{1,3\}\)
has an irreducible first block.
\item[(A)] Every absolutely irreducible characteristic-zero point
of \(S\) has the cuspidal attachment over the fixed quadratic
field selected in Theorem~\ref{thm:rh-symplectic-spectrum},
with a common choice of coefficient isomorphism.
\end{enumerate}
These are not additional hypotheses on top of the original
\texttt{Fin0}, \texttt{NoEnd}, \texttt{AllIrr}:
the first is equivalent to \texttt{Fin0} on this universal
model, the second is the literal endpoint condition, and the
third is weaker than the original automorphy premise.
The proof below establishes every remaining passage from these
premises to the individual uniform integral bound. It does
not assume a bound on the lengths or a bound on the number of
integral equations.

\begin{lemma}[Integral coordinate projectors and the support graph]
\label{lem:rh-sp4-graph}
Let \(u\) be an actual characteristic-zero point. Let
\(i\rightsquigarrow j\) mean that some matrix \(u(g)\) has
nonzero \(ij\)-entry. The relation is reflexive, transitive,
and is preserved by \((i,j)\mapsto(3-j,3-i)\).
Exactly the following alternatives exhaust the possibilities
(up to overlap of their descriptions):
\begin{enumerate}
\item the relation is full, and \(u\) is absolutely irreducible;
\item every mutual pair is diagonal;
\item both prescribed coordinate blocks are strongly connected,
and at least one direction between these blocks is entirely absent.
\end{enumerate}
In the second case every nonzero matrix entry has
\(\mu_i-\mu_j\ge0\), or every such entry has
\(\mu_i-\mu_j\le0\). The alternative is determined by
\(u(s_0)_{00}\), which is \(\tau(s_0)\) in the first ordering
and \(\alpha(s_0)\) in the reversed ordering.
\end{lemma}

\begin{proof}
By the finite-coefficient lemma in
Section~\ref{sec:rh-symplectic-point-rigidity}, we work over
the integers \(O'\) of a finite coefficient extension, in the
given normalized basis. Its fraction field can be enlarged
when necessary. Write \(D_1,D_2\) for the two coordinate
blocks. The block-normalization of \(u(\mathrm{ph})\) and
the residual separation (BU8) give relatively prime
characteristic polynomials of its two blocks over \(O'\):
their resultant is a unit. The Bezout identity and
Cayley--Hamilton therefore put the two block projectors in
the \(O'\)-algebra generated by the matrices \(u(g)\).

Write \(H=u(g_0)+\tau(g_0)u(g_0)^{-1}\).
For \(X^\dagger=J^{-1}X^{\mathsf t}J\), with the alternating
form \(J\) in Theorem~\ref{thm:rh-symplectic-local},
we have \(H^\dagger=H\). Its two compressed block matrices
are diagonal by the two marking equations (SM12) and their
reflections. In the first block the two diagonal entries
differ modulo the maximal ideal by
\[
          (\bar\alpha(g_0)-1)(\bar\Psi(g_0)-1).
\]
The second block has the opposite ordering of the same two
residual entries. This difference is a unit by the actual
choice of \(g_0\). Taking the linear spectral projector of
each compressed block now puts every coordinate matrix
unit \(E_{ii}\) in the same integral algebra.

Consequently the algebra over the coefficient field
generated by \(u(g)\) contains \(E_{ij}\) whenever
\(i\rightsquigarrow j\); conversely its support is contained
in the union of these supports since it is the linear span
of the group matrices. Multiplication of matrix units gives
transitivity. The fixed similitude gives
\(u(g)^\dagger=\tau(g)u(g)^{-1}\), so reflection preserves
the relation. Full support is equivalent to absolute
irreducibility by the matrix-algebra criterion.

\emph{Excluding a two-element component across the blocks.}
Suppose distinct \(i,j\) in different blocks lie in one
strongly connected component of size two. The coordinate
span of all indices reachable from \(i\) is invariant.
Its coordinate subspace obtained by deleting \(i,j\) is
also invariant: an arrow back to \(i\) or \(j\) would
enlarge the component. The two-dimensional subquotient
has matrices with the four entries \(ii,ij,ji,jj\) of \(u\).
Both off-diagonal entries occur nontrivially. The ordinary
local flag, restricted to this subquotient, gives a
local invariant line. Since \(\mathrm{ph}\) is local,
is block diagonal, and has distinct eigenvalues on
these two coordinates, this line must be a coordinate
line. Relabel \(i,j\), if necessary, so that the \(ij\)-entry
vanishes on the decomposition group \(D_\ell\).

Choose \(b\in O'\) of minimal valuation among the nonzero
values of this \(ij\)-entry and put
\[
 c(g)=u(g)_{ij}/b,\qquad
 z(g)=\bar c(g)/\bar\chi_j(g),
 \qquad \chi=(\tau,\Psi,\alpha,1).
\]
The two-dimensional multiplication law gives
\(z(gh)=z(g)+(\bar\chi_i/\bar\chi_j)(g)z(h)\),
and \(z\ne0\) by minimality. It is continuous with values
in the residue field, is zero on \(D_\ell\), and satisfies
the original inertia-kernel vanishing at all other finite
places. The residual quotient character is one of the
six off-block roots \(\alpha,\beta,\tau\) or their inverses.
The actual \texttt{XVan} proved in
Theorem~\ref{thm:rh-symplectic-arithmetic-data} makes this
cocycle a global coboundary. This assertion persists under
finite extension of the residue field, by decomposing the
cocycle in a basis over the original residue field.
At \(\mathrm{ph}\) the residual quotient character differs
from one, so the primitive of that coboundary is zero.
This contradicts \(z\ne0\). In this singleton-local case
the block normalization thus removes the need for a
separate estimate on an invariant-line slope.

A strongly connected component of size three would be
sent by reflection to another such component, necessarily
itself. This is impossible because the reversing involution
has no fixed index. Hence a mutual cross-block pair forces
full support. If there are no such pairs but there is
a nondiagonal mutual pair, reflection gives exactly the
two strongly connected coordinate blocks. Transitivity
forbids arrows in both directions between them. This
proves the three alternatives.

\emph{The ordering in the diagonal case.}
In the absence of nondiagonal mutual pairs the support is
a partial order. The four diagonal entries are global
characters, denoted \(\kappa_i\), by matrix multiplication;
the representation is triangular for any linear
extension of that partial order. The local flag has the
four distinct prescribed inertia characters, so the
diagonal local characters agree with those of that flag
in a permutation \(m_1\) of the four indices (compare
the two sums of characters, using linear independence
of distinct characters). By (SR1) the local graded
values at \(\mathrm{ph}\) are the original labeled
values. Their block separation forces
\(\lambda_{m_1(i)}=\lambda_i\).
The fixed similitude and triangularity give
\(\kappa_i\kappa_{3-i}=\tau\). The distinct eigenvalues
at \(s_0\) and their prescribed partners then imply
\(m_1(3-i)=3-m_1(i)\). There are only two such
permutations, the identity and \((0,1,2,3)\mapsto(2,3,0,1)\).
In particular
\[
         \kappa_i|_{I_\ell}=\chi_{m_1(i)}|_{I_\ell}.
\]
Writing \(h=(e+1,e,1,0)\), each \(\kappa_i\) is
\(\varepsilon^{h_{m_1(i)}}\) times a finite-order character.
Indeed this ratio has finite inertia image everywhere
and is unramified outside a finite set. A common power
is unramified at every finite prime, and class field
theory for \(\mathbb Q\) makes that power trivial.

For each cover \(i\rightsquigarrow j\) in the partial order
there is a two-dimensional coordinate subquotient:
take all coordinates reachable from \(j\) and delete
\(i,j\), which is invariant by the absence of an
intermediate arrow. It is an extension of \(\kappa_j\)
by \(\kappa_i\). If \(m_1(i)>m_1(j)\), the inherited local
ordinary flag on this subquotient begins with \(\kappa_j\).
This follows from the distinct spectral values at \(s_0\):
the original flag is the sequence of sums of its ordered
eigenlines, and restriction and quotient by invariant
subspaces preserve the surviving order. Thus the extension
splits at \(\ell\). Away from \(\ell\) it is a local
inertia coboundary, since the representation is trivial
on a finite-index normal subgroup of inertia and
characteristic-zero finite-group first cohomology is zero.
Its character \(\kappa_i/\kappa_j\) is a strictly negative
cyclotomic power times a finite-order character.
Proposition~\ref{prop:selmer-negative-weight-rigidity}
therefore splits the extension globally. It would follow
that \(u(g)_{ij}=a(\kappa_i(g)-\kappa_j(g))\) for a fixed
scalar \(a\), for every \(g\).
The same linear equality would hold on the image algebra.
Evaluate it on \(E_{ii}\), which belongs to this algebra:
it forces \(a=0\), contradicting the cover.
Hence every cover, and by transitivity every strict
arrow, satisfies \(m_1(i)<m_1(j)\).
For the identity this implies \(\mu_i\ge\mu_j\);
for the other ordering it implies \(\mu_i\le\mu_j\).
Their \(00\)-entries at \(s_0\) are the stated distinct
values. No bound on the coefficients of a coboundary
is used.
\end{proof}

\begin{lemma}[The actual one-torus slice and its rational chart]
\label{lem:rh-sp4-slice}
Let
\[
 \mathscr T=\prod_{j} R_{(0,j)},\quad
 A_1=\bigoplus_j R_{(0,j)},\quad B_1=R_{(0,0)},\quad
 Q_1=(F,E_+,Z)A_1.
\]
The product notation means the subring of first-weight-zero
formal series in the complete ring \(R\), not a Cartesian
product of rings. Then \(B_1\) is complete Noetherian local,
\(A_1\) is Noetherian and generated over \(B_1\) by finitely
many elements of second weights \(-2,-1,1,2\), and
\[
                 A_1/Q_1=\bigoplus_j S_{(0,j)}.           \tag{SL1}
\]
The rational augmentation cotangent space of \(A_1\) has
multiplicities
\[
 \begin{array}{c|ccccc}
 j&-2&-1&0&1&2\\ \hline
 \dim&0&r&m&r'+1&1 .
 \end{array}                                               \tag{SL2}
\]
There is a chart inducing this weightwise identification,
with the original augmentation as its constant coefficient.
The negative and zero cotangent spaces are spanned over
the fraction field by the corresponding original equations.
\end{lemma}

\begin{proof}
Take the homogeneous power-series presentation of the
actual ambient ring \(R\). The finitely many variables
are entry displacements; their weights have both
coordinates in \(\{-1,0,1\}\). Weight projection on
formal series descends through the defining ideal,
which is stable under every weight projection.
All first-weight-zero monomials factor into first-zero
variables and pairs of variables of first weights \(1,-1\).
This proves that \(\mathscr T\) is a quotient of a
power-series ring on these finitely many monomials.
Its presentation is equivariant for the second grading.
The weight-zero subalgebra of such a ring is itself a
quotient of a finite-variable power-series ring: the
semigroup of nonnegative exponent vectors of weight zero
is finitely generated by its nonzero componentwise
minimal elements (Dickson's lemma and induction on
total degree). The same minimal-element argument for
each prescribed weight expresses a homogeneous series
as a finite sum of monomials of that weight times
weight-zero series. The generators of the finite-sum
algebra may thus be taken to be the original first-zero
variables and the paired monomials. Their second weights
belong to \(\{-2,-1,0,1,2\}\). The inverse of a weight-zero
series with unit constant coefficient is weight zero,
so \(B_1\) is local. All these statements concern the
full integral rings. Since the equation generators
are homogeneous of first weight zero, projecting their
coefficients to the required weights proves (SL1)
and in particular
\(B_1/(Q_1\cap B_1)=S_{(0,0)}\).

For the chart, let \(\mathfrak a\) and \(\mathfrak b\)
be the augmentation ideals in \(R\) and \(\mathscr T\).
Their filtrations on the first-invariant ring are
cofinal:
\[
 \mathfrak a^{2n}\cap\mathscr T\subseteq\mathfrak b^n
 \subseteq\mathfrak a^n\cap\mathscr T .
\]
For the first inclusion lift to series of order at
least \(2n\) and project onto first weight zero;
every such monomial factors into at least \(n\)
of the generators of degrees at most two.
Grouping the resulting series by finitely many
prefixes of length \(n\) gives membership in
\(\mathfrak b^n\). The projection also shows that
the invariant part of each finite augmentation jet
is the quotient by the intersected ideal, and it
commutes with inverting \(O\setminus\{0\}\).
These jets are finite \(O\)-modules.
On the slice the augmentation ideal is generated by
the same homogeneous generators. Its powers intersect
\(\mathscr T\) and \(A_1\) compatibly: an expression
in the generators of such a power can be projected
onto the finitely many second weights of an element
of \(A_1\). Moreover the truncated monomial expansions
give surjectivity from \(A_1\) onto every
\(\mathscr T/\mathfrak b^n\).

It follows that the rational augmentation completion
of \(A_1\) is precisely the first-weight-zero ring
of the rational augmentation chart of \(R\).
In the latter chart there are the variables of
first weight zero and exactly the three additional
variables \(u_A,u_B,u_T\).
Every first-invariant monomial in these three is a
unique product of \(u_Au_B\) and \(u_Tu_B\), with
weights \(P\) and \(2P\).
Thus the invariant chart is a power-series ring on
the original first-zero variables and these two
additional independent coordinates. This proves (SL2).
Its linear coefficient map identifies the rational
cotangent space of \(A_1\), since the augmentation
jets just identified are the actual jets after
inversion of the uniformizer.

It remains to check spanning by the equations.
At weights \(-1,0\) the slice cotangent space
has exactly the original ambient coordinates of
weights \(-P,0\), since both added variables have
positive second weight. A rational linear functional
on either space that kills the corresponding equations
extends to an augmentation derivation on \(R\)
in that full weight, killing the other equations
by homogeneity. It therefore gives an actual normalized
dual-number point of \(S\); the finite augmentation
jet ensures bounded integral denominators for its
coefficient realization.
Proposition~\ref{prop:rh-symplectic-local-finite}
puts its cocycle in the finite local conditions.
At weight \(-P\), its adjoint character is
\(\Psi^{-1}\), and
Proposition~\ref{prop:selmer-negative-weight-rigidity}
gives vanishing. At weight zero the adjoint
characters are trivial, so finite cohomology
vanishes by class field theory for \(\mathbb Q\).
Thus the cocycle is a global coboundary.

For clarity, the marking kills all nontrivial
coboundary tangents at the base. Any primitive
\(X\in\mathfrak{sp}_4\) of a normalized
coboundary must be block diagonal, by differentiating
the block condition at \(\mathrm{ph}\) and using
the disjoint block spectra. Differentiating the two
marked entries of \(H\) then forces the two off-diagonal
entries of the first block of \(X\) to be zero,
since its compressed \(H\) is diagonal with
unit difference; the symplectic relation gives the
other block as well. Hence \(X\) is diagonal.
It commutes with \(r_0\), making the base tangent
zero. By the generation of the universal ring by
its actual entries this also makes the derivation zero.
There was no nonzero functional annihilating the
equations. Rational spanning at \(-2\) is automatic
since its cotangent space has dimension zero.
\end{proof}

\begin{lemma}[Lifting slice points and tangents]
\label{lem:rh-sp4-slice-points}
Assume (F) and put \(\mathscr S=\bigoplus_{i,j}S_{(i,j)}\).
Every \(\overline E_0\)-point of \(A_1/Q_1\) lifts to a
point of \(\mathscr S\), which becomes an actual local
characteristic-zero point of \(S\) after rescaling by the
two-dimensional torus. If the lifted representation is
absolutely irreducible, every slice derivation at the given
point lifts to an actual normalized dual-number tangent,
after the same rescaling. Under the rigidity implication
of Theorem~\ref{thm:rh-symplectic-rigidity}, the slice
tangent space has dimension at most one.
Moreover isomorphic actual irreducible representations
arising from two slice points imply that those slice
points belong to one second-torus orbit.
\end{lemma}

\begin{proof}
First, (F) follows from the original \texttt{Fin0}
on the universal representation and is equivalent to it
on this model. The integral algebra of balanced
entry monomials is dense in \(S_{(0,0)}\):
the entries topologically generate \(S\), so approximate
by polynomials and project the polynomial to weight
\((0,0)\). The powers of the maximal ideal are stable
under weight projection. If this algebra is finite
over \(O\), it is compact and hence closed in the
complete local ring, giving equality with \(S_{(0,0)}\).
Conversely finiteness of \(S_{(0,0)}\) makes the
submodule of balanced entry monomials finite, and
every local specialization has the same conclusion.

Projection onto first weight zero is a linear
retraction of \(\mathscr S\) over \(A_1/Q_1\).
Consequently the fibre algebra of \(\mathscr S\) over
a given slice point is nonzero. It is a finite-type
algebra over the algebraically closed field:
the homogeneous monomial factorization by Dickson's
lemma expresses \(\mathscr S\) as generated by its
finitely many entry-displacement generators over
the finite \(O\)-algebra \(S_{(0,0)}\).
The weak Nullstellensatz gives a lifted point.

We justify the rescaling with all invariant series retained.
For each generator with nonzero value \(a_i\), set
\(v_i=v_\ell(a_i)\), and write \(w_i\in\mathbb Z^2\)
for its weight in this paragraph. Every nonzero
nonnegative integral relation
\(\sum_i b_i w_i=0\) satisfies
\(\sum_i b_i v_i>0\). The corresponding invariant
monomial belongs to the maximal ideal of the finite
local ring \(S_{(0,0)}\), all of whose characters
have positive valuation on that ideal.
There exists a rational vector \(s\in\mathbb Q^2\)
such that
\[
                         v_i+w_i\cdot s>0
\]
for every nonzero generator value.
Indeed the polytope \(b_i\ge0,\ \sum b_i=1,\
\sum b_iw_i=0\) has rational vertices.
If nonempty, \(\sum b_i v_i\) is strictly positive
on every vertex by clearing denominators.
Choose a positive rational \(\eta\) less than its
minimum. The finite-dimensional linear inequality
alternative (Farkas' lemma) gives feasibility of
\(w_i\cdot s\ge\eta-v_i\), since every nonnegative
linear relation among the \(w_i\)'s has nonpositive
dot product with the vector \((\eta-v_i)_i\).
An empty polytope has no such nonzero obstruction.
Rational solutions exist as the constraints are
rational. Choose torus parameters of valuations
\(s_1,s_2\) in \(\overline E_0^\times\).

After rescaling, all entry-displacement generators
have strictly positive valuation, including the zero
values by convention. Values of the finite algebra
\(S_{(0,0)}\) and of the finitely many generators lie
in one finite coefficient extension; the entire
graded algebra has integral values. Substitution
therefore defines a continuous map on the power-series
presentation. To verify agreement with the lifted
point on each homogeneous series, group the tail
of order at least \(n\) by its finitely many
degree-\(n\) monomial prefixes. Each remaining
factor is a homogeneous series, with integral
value by the finite-generation argument.
Thus the evaluated tail tends to zero.
The same argument on every homogeneous component
of a relation proves that the continuous map kills
the entire defining ideal, not just a characteristic-zero
quotient. This is the required local point of \(S\).
Its effect on the slice is just rescaling by the
second torus.

Now take an irreducible lifted representation.
There must be some nonzero entry \(a\) of first
degree \(1\) and some nonzero entry \(b\) of first
degree \(-1\), otherwise a coordinate block is
invariant. Put \(f=ab\), in \(A_1/Q_1\).
At the lifted point \(f\ne0\).
The first grading gives an exact algebra isomorphism
\[
       \mathscr S[1/f]=(A_1/Q_1)[1/f][a,a^{-1}].
\]
In fact \(a\) is an invertible element of first degree
one after localization; multiplying any homogeneous
element by an appropriate power of \(a\) puts it
in first degree zero. Independence of the grading
also proves the displayed Laurent-polynomial
description with no extra relation.
Extend the slice dual-number map by prescribing
the scalar value of \(a\). It is a map on \(\mathscr S\).
After torus transport its scalar part is the
local point already obtained. Its image lies in
\[
 O''+\delta\varpi^{-b_0}O'',\qquad \delta^2=0,
\]
for some finite coefficient DVR \(O''\) of uniformizer
\(\varpi\) and some \(b_0\ge0\): the finitely many
generators and the finite \(O\)-module \(S_{(0,0)}\)
supply all bounds. Every product of \(n\) generator
values has at most one nilpotent contribution, so
it tends uniformly to zero as \(n\) grows.
The same homogeneous-tail argument extends the
map to \(S\) and verifies every relation.
This realizes the tangent in the actual local model.

At a rigid irreducible point its tangent cocycle
has a primitive \(X\in\mathfrak{sp}_4\). Exactly the
block and marked differentiation in
Lemma~\ref{lem:rh-sp4-slice} forces \(X\) to be
diagonal, with entries \((a_0,b_1,-b_1,-a_0)\).
Modulo scalar matrices this is the infinitesimal
action of the two original torus parameters.
The first acts trivially on the slice, so the
restriction of every actual tangent belongs to
the single second-torus direction.
The actual representation determines its dual-number
map by generation, hence this is also true for
the lifted algebraic tangent. This proves the
one-dimensional bound.

Finally suppose two actual irreducible points are
conjugate by \(U\in\mathrm{GL}_4(\overline E_0)\).
The spectral sets of \(\mathrm{ph}\) on the two
blocks are fixed, by (SR1), and disjoint.
Thus \(U\) preserves the coordinate blocks.
The compressed marked matrices \(H\) have diagonal
entries with two distinct labeled residual values.
The possible matching of their eigenvalues between
the two points is necessarily label-preserving,
because the two residue neighborhoods are disjoint.
The intertwiner \(U\) is therefore diagonal.
The fixed similitude at both points and Schur's
lemma imply \(U^{\mathsf t}JU\) is a scalar multiple
of \(J\). Thus its diagonal entries satisfy
\(u_0u_3=u_1u_2\). Up to scalar the matrix is exactly
\(\operatorname{diag}(t_1t_2,t_2,t_1,1)\).
The continuous maps determined by the conjugate
representations agree with that torus transformation
on every homogeneous series, by generation and
continuity. Restriction to the slice gives just its
second-torus orbit; the preliminary rescalings also
have this property.
\end{proof}

\begin{theorem}[Conditional uniform bound for the actual RH Case II length]
\label{thm:rh-sp4-conditional-length}
Assume precisely (F), (E), and (A) for the actual datum.
Then the original strict character length of
Proposition~\ref{shared-rh-integral-product} satisfies
\[
 (\log\ell)\ell(\ell,e,\psi)
               \le(\ell N(e+1))^{C},                       \tag{SL3}
\]
with an absolute exponent \(C\), independent of the
coefficient field, the number of generators, and the
number and depths of the original relations.
In particular no additional rank-four integral counting
or compatibility hypothesis is needed after those three
point-arithmetic premises.
\end{theorem}

\begin{proof}
We check the point hypotheses of the abstract length theorem
\ref{thm:selmer-abstract-length}, with the exact slice of
Lemma~\ref{lem:rh-sp4-slice} and \(s=s'=0\).
The finite degree-zero quotient is (F).
For every lifted slice point the vanishing of a whole
sign of the second weights in the representation implies
vanishing of that sign on the slice: write homogeneous
series as finite sums of monomials in the entry
displacements with coefficients of full weight zero.
A nonzero monomial of negative second weight has
a factor of negative second weight, and similarly
for positive second weight. Diagonal displacements
have second weight zero. This argument works on the
actual local specialization and, by equivariance,
before rescaling as well.

If a two-sided slice point lifted to a representation
of the diagonal alternative of Lemma~\ref{lem:rh-sp4-graph},
one sign would vanish, a contradiction.
If it lifted to the third alternative, one direction
between the two coordinate blocks would be absent.
Conjugate by first-torus parameters tending to zero or
to infinity, in the direction which multiplies the
remaining off-block entries by a power of a uniformizer.
Their original diagonal block entries do not change.
In the complete ring, entry-displacement generators have
positive valuation; the rescaled values still do,
and converge to the values obtained by setting the
off-block generators to zero. Continuous power-series
evaluation gives an actual limiting point of \(S\),
since every original relation vanishes before the limit.
It is block diagonal with the same irreducible first
block, forbidden by (E). Hence all two-sided points
lift to irreducible representations.

For fixed-point separation, start with a slice
point killing all nonzero second weights.
Its lifted representation has the diagonal alternative:
in each of the other two alternatives both entries
\(02\) and \(20\) occur nontrivially and have opposite
nonzero slice weights. Any other point with the same
degree-zero character also has the diagonal alternative,
because every product of opposite matrix entries
\(u(g)_{ij}u(h)_{ji}\) is of full weight zero.
These products vanish for the first representation
and hence for the second. Moreover their \(00\)-entries
at \(s_0\) agree, being degree zero.
Lemma~\ref{lem:rh-sp4-graph} now forces the same
ordering and hence the same vanishing sign throughout
that fibre. The abstract fixed-point separation
condition is proved, even without its exemption
of the augmentation point.

By (A), Theorems~\ref{thm:rh-symplectic-spectrum}
and \ref{thm:rh-symplectic-rigidity} give rigidity
and a common nonzero trace difference for every
irreducible lift. Lemma~\ref{lem:rh-sp4-slice-points}
then gives the one-dimensional slice tangent bound.
Take the trace element \(a\in B_1\) to be any
homogeneous lift of
\(\operatorname{tr}\rho_u(\sigma)-\operatorname{tr}r_0(\sigma)\).
Its augmentation is zero. By (FS4) it is nonzero
at every two-sided slice point.
The exact chart, equation counts, spanning and
generation hypotheses are those of
Lemma~\ref{lem:rh-sp4-slice}. This verifies every
hypothesis of the abstract length theorem.

Let \(I=\ker(A_1\to O)\),
\(I^{(2)}\) the saturation of \(I^2\) by nonzero
scalars of \(O\), and let
\(\mathcal C^-\) be the quotient of
\((I/I^{(2)})_{-1}\) by the classes of the \(r+1\)
actual negative equations.
Take one representative of every two-sided slice
orbit, with a corresponding actual irreducible
representation. By the last part of
Lemma~\ref{lem:rh-sp4-slice-points}, the chosen
representations are pairwise nonisomorphic.
Denote the resulting finite family by \(X\).
The abstract length theorem gives
\[
   \operatorname{length}_O\mathcal C^-
       \le2\sum_{u\in X}
        v_\varpi\big(\operatorname{tr}u(\sigma)
                          -\operatorname{tr}r_0(\sigma)\big),
                                                               \tag{SL4}
\]
where \(\varpi\) is a uniformizer of \(O\).
The valuations are those at the actual points; no
intersection of ideals with a reduced special fibre
has been substituted for the original integral cotangent.

We prove the lower comparison for the whole strict
Selmer group, not just for its rationalized tangent.
Put \(\theta=\Psi^{-1}\) and \(D=(E_0/O)(\theta)\).
The strict group in Proposition~\ref{shared-rh-integral-product}
is finite by the qualitative finiteness theorem there
(or by prime-to-\(\ell\) restriction to the host and
Proposition~\ref{shared-selmer-finite}).
Let \(b\ne0\) in \(O\) annihilate it.
Every class has a unique cocycle representative
\(z\) with \(z(g_0)=0\), obtained by subtracting the
coboundary with primitive \(z(g_0)/(\theta(g_0)-1)\).
This is well defined because the denominator is a unit.
The assignment is \(O\)-linear. In particular
\(bz(g)=0\) everywhere.
The restriction to \(D_\ell\) is a coboundary with
primitive \(m_\ell\) in \(D\).
Choose \(s_0\) as in the original local datum;
\(d_0=\theta(s_0)-1\ne0\).
Since \(b m_\ell\) is locally invariant,
\(b d_0m_\ell=0\). Set \(c_0=bd_0\ne0\).
Identify \(O/(c_0)\) with \(D[c_0]\) by multiplication
by \(1/c_0\). All normalized cocycles then have values
in this fixed finite module and are coboundaries
on \(D_\ell\) in the very same module.
Away from \(\ell\), the strict local condition makes
them zero on the inertia kernel of the residual
diagonal representation: there the prime-to-\(\ell\)
finite inertia characters have the same kernels
before and after reduction.
These are exactly the premises of the finite-modulus
derivation of Theorem~\ref{thm:rh-symplectic-model}.
The enlargement from \(b\) to \(c_0\) is important:
the finite-order local primitive need not have been
killed by \(b\) itself.

For each class, that theorem gives an integral
derivation on \(R\), of weight \(-P\), which modulo
\(c_0\) kills the full equation ideal. Its restriction
to \(A_1\) kills \(I^2\) integrally; as its target
before reduction is torsion-free, it also kills
\(I^{(2)}\).
It therefore induces a homomorphism from
\(\mathcal C^-\) to \(O/(c_0)\).
This homomorphism is independent of the chosen
integral lift and is linear in the Selmer class:
the modular derivation on \(S\) is determined by
the matrices (SM13), and their nilpotent terms
are linear in the normalized cocycle.
The resulting map from the strict group into
\(\operatorname{Hom}_O(\mathcal C^-,O/(c_0))\)
is injective. Indeed the derivative of a homogeneous
lift of the \(20\)-entry at every \(g\) equals
\(\tau(g)z(g)\) modulo \(c_0\), by (SM14).
A zero homomorphism makes every \(z(g)\) zero.
Smith normal form gives
\(\operatorname{length}_O\operatorname{Hom}_O
(\mathcal C^-,O/(c_0))\le
\operatorname{length}_O\mathcal C^-\).
This establishes the integral lower comparison for
the entire group. It is the rank-four use of the
original modular derivation field, with no uniform
annihilator assumed.

Divide (SL4) by the ramification index of \(O\)
and use \(v_\varpi=b_Ov_\ell\).
The integral Fourier normalization in
Proposition~\ref{shared-rh-integral-product}
then gives
\[
 (\log\ell)\ell(\ell,e,\psi)
 \le2(\log\ell)\sum_{u\in X}
    v_\ell\big(\operatorname{tr}u(\sigma)
                            -\operatorname{tr}r_0(\sigma)\big).
\]
The fixed-field estimate (FS4) proves (SL3).
All constants in it are independent of the size
of the presentation, as required.
\end{proof}

\paragraph{Scope and the remaining arithmetic.}
If (F), (E), (A) are supplied uniformly for the
remaining actual RH Case II data, (SL3) bounds every
term left in (DC11). By (RN4) all these primes are
at most \(2e+1\), and the number of characters is
at most \(N\). Summing gives (RN7) for the remaining
resonant twists. The already proved nonresonant and
exponent-one estimates and the original supplied
remainders of the RH product must still be applied
on the same windows. This theorem is a conditional
comparison on the exact model, not a proof of (F),
(E), or (A). It does not replace those arithmetic
premises by generic smoothness, semisimplification,
or a bound on the reduced host degree.
